\documentclass[11pt, a4paper]{article}

% ============================================================================
% PACKAGES
% ============================================================================
\usepackage[utf8]{inputenc}
\usepackage[T1]{fontenc}
\usepackage{amsmath, amssymb, amsthm}
\usepackage{mathtools}
\usepackage{bm}              % Bold math symbols
\usepackage{geometry}
\usepackage{booktabs}
\usepackage{array}
\usepackage{xcolor}
\usepackage{tcolorbox}
\usepackage{tikz}
\usepackage{hyperref}
\usepackage{enumitem}

\geometry{margin=2.5cm}

% ============================================================================
% CUSTOM ENVIRONMENTS
% ============================================================================

% Formal definition box
\newtcolorbox{formaldef}[1][]{
  colback=blue!5!white,
  colframe=blue!60!black,
  fonttitle=\bfseries,
  title={#1},
  sharp corners,
  boxrule=0.8pt
}

% Intuitive explanation box
\newtcolorbox{intuition}[1][]{
  colback=green!5!white,
  colframe=green!50!black,
  fonttitle=\bfseries,
  title={#1},
  sharp corners,
  boxrule=0.8pt
}

% Transformer application box
\newtcolorbox{transformer}[1][]{
  colback=orange!5!white,
  colframe=orange!60!black,
  fonttitle=\bfseries,
  title={#1},
  sharp corners,
  boxrule=0.8pt
}

% Dimension check box
\newtcolorbox{dimcheck}{
  colback=gray!8!white,
  colframe=gray!50!black,
  fonttitle=\bfseries,
  title={Dimension Check},
  sharp corners,
  boxrule=0.6pt
}

% ============================================================================
% NOTATION
% ============================================================================
\newcommand{\R}{\mathbb{R}}
\newcommand{\vx}{\bm{x}}
\newcommand{\vq}{\bm{q}}
\newcommand{\vk}{\bm{k}}
\newcommand{\vv}{\bm{v}}
\newcommand{\ve}{\bm{e}}
\newcommand{\vp}{\bm{p}}
\newcommand{\vo}{\bm{o}}
\newcommand{\mW}{\bm{W}}
\newcommand{\mQ}{\bm{Q}}
\newcommand{\mK}{\bm{K}}
\newcommand{\mV}{\bm{V}}
\newcommand{\mX}{\bm{X}}
\newcommand{\mE}{\bm{E}}
\newcommand{\mP}{\bm{P}}
\newcommand{\mA}{\bm{A}}
\newcommand{\mO}{\bm{O}}
\newcommand{\mI}{\bm{I}}
\newcommand{\shape}[1]{\underset{\scriptscriptstyle #1}{\vphantom{X}}}

% ============================================================================
% DOCUMENT
% ============================================================================

\title{\textbf{The Linear Algebra of Attention} \\[0.5em]
\large A Complete Mathematical Reference for the Proust Attention Machine}
\author{Proust Attention Project}
\date{January 2026}

\begin{document}
\maketitle

\begin{abstract}
This document maps every linear algebra operation in our transformer architecture,
from character embedding to attention output. Each concept is presented in three layers:
\textbf{formal definition} (blue box), \textbf{geometric intuition} (green box),
and \textbf{transformer application} (orange box). Read it to connect your knowledge
of linear algebra to what the Proust machine actually computes.
\end{abstract}

\tableofcontents
\newpage

% ============================================================================
\section{Notation and Conventions}
% ============================================================================

\begin{formaldef}[Notation Table]
\begin{tabular}{@{}lll@{}}
\toprule
\textbf{Symbol} & \textbf{Meaning} & \textbf{Value in our model} \\
\midrule
$V$ & Vocabulary size (number of unique characters) & $\approx 80$ \\
$d$ & Model dimension (\texttt{d\_model}) & $64$ \\
$d_k$ & Dimension per attention head & $32$ \\
$n_h$ & Number of attention heads & $2$ \\
$T$ & Sequence length (\texttt{seq\_len}) & $\leq 256$ \\
$B$ & Batch size & variable \\
$\R^{m \times n}$ & Set of all real matrices with $m$ rows and $n$ columns & --- \\
$\bm{A}^\top$ & Transpose of matrix $\bm{A}$ & --- \\
$\bm{A} \bm{B}$ & Matrix multiplication of $\bm{A}$ and $\bm{B}$ & --- \\
$\langle \vx, \bm{y} \rangle$ & Dot product (inner product) of vectors $\vx$ and $\bm{y}$ & --- \\
\bottomrule
\end{tabular}
\end{formaldef}

\begin{intuition}[How to Read Shapes]
Throughout this document, we annotate shapes as $(B, T, d)$ meaning:
\begin{itemize}[nosep]
  \item First axis: which sequence in the batch ($B$ sequences)
  \item Second axis: which position in the sequence ($T$ positions)
  \item Third axis: which feature/dimension ($d$ numbers per position)
\end{itemize}
A matrix multiplication $(B, T, \mathbf{m}) \times (\mathbf{m}, n) = (B, T, n)$ requires
the \textbf{inner dimensions} to match (both are $m$).
\end{intuition}

% ============================================================================
\section{Vectors and Vector Spaces}
% ============================================================================

\begin{formaldef}[Vector in $\R^n$]
A \textbf{vector} $\vx \in \R^n$ is an ordered list of $n$ real numbers:
\[
\vx = \begin{pmatrix} x_1 \\ x_2 \\ \vdots \\ x_n \end{pmatrix} \in \R^n
\]
The set $\R^n$ forms a \textbf{vector space}: vectors can be added and scaled,
and the results remain in $\R^n$.
\end{formaldef}

\begin{intuition}[Vectors as Points in Space]
A vector in $\R^{64}$ is a point in 64-dimensional space. We can't visualize 64 dimensions,
but the math works identically to 2D or 3D:
\begin{itemize}[nosep]
  \item Two vectors can be \textbf{close} (similar) or \textbf{far} (different)
  \item They can \textbf{point in the same direction} (aligned) or \textbf{perpendicular} (unrelated)
  \item We can \textbf{add} them (combine information) and \textbf{scale} them (amplify/diminish)
\end{itemize}
Every character in our model is represented as a point in $\R^{64}$.
\end{intuition}

\begin{transformer}[Embedding Vectors]
Each character in the vocabulary gets a vector $\ve_i \in \R^{64}$.
After training, characters that behave similarly in Proust's text will have
vectors that are close together in this 64-dimensional space.
Vowels might cluster. Consonants might cluster. Punctuation in another region.
At initialization, these vectors are random --- the structure emerges from training.
\end{transformer}

% ============================================================================
\section{The Embedding Lookup: Indexing Into a Matrix}
% ============================================================================

\begin{formaldef}[Embedding as Row Selection]
Let $\mE \in \R^{V \times d}$ be the \textbf{embedding matrix}, where row $i$ is the
embedding vector for character $i$ in the vocabulary.

Given a sequence of token IDs $(t_1, t_2, \ldots, t_T)$ where each $t_j \in \{0, 1, \ldots, V-1\}$,
the embedding operation selects rows:
\[
\mX_{\text{tok}} = \begin{pmatrix} \ve_{t_1}^\top \\ \ve_{t_2}^\top \\ \vdots \\ \ve_{t_T}^\top \end{pmatrix}
\in \R^{T \times d}
\]
This is equivalent to multiplication by a one-hot matrix:
$\mX_{\text{tok}} = \bm{H} \mE$ where $\bm{H} \in \{0,1\}^{T \times V}$
and $H_{j,i} = 1$ iff $t_j = i$.
\end{formaldef}

\begin{dimcheck}
\[
\underbrace{\bm{H}}_{T \times V} \times \underbrace{\mE}_{V \times d} = \underbrace{\mX_{\text{tok}}}_{T \times d}
\]
Inner dimension $V$ is consumed. Each of $T$ positions gets a $d$-dimensional vector.
\end{dimcheck}

\begin{intuition}[Lookup vs.\ Multiplication]
Selecting row 5 from a matrix is the same as multiplying by a one-hot vector
$[0,0,0,0,0,1,0,\ldots]$. In code, we use direct indexing (\texttt{E[token\_ids]})
because it's faster than constructing a one-hot matrix. But mathematically, embedding
\textit{is} a matrix multiplication --- which is why it's differentiable and learnable in gradient terms.
\end{intuition}

\begin{transformer}[In the Proust Machine]
\texttt{token\_embedding[token\_ids]} on line 95 of \texttt{embedding.py}:
\begin{itemize}[nosep]
  \item Input: $(B, T)$ integers --- each integer is a character ID
  \item Table: $(80, 64)$ --- one row of 64 numbers per character
  \item Output: $(B, T, 64)$ --- each character is now a 64-dim vector
\end{itemize}
\end{transformer}

% ============================================================================
\section{Vector Addition: Combining Embeddings with Position}
% ============================================================================

\begin{formaldef}[Vector Addition]
For vectors $\vx, \bm{y} \in \R^n$, their sum is defined element-wise:
\[
\vx + \bm{y} = \begin{pmatrix} x_1 + y_1 \\ x_2 + y_2 \\ \vdots \\ x_n + y_n \end{pmatrix} \in \R^n
\]
The result lies in the same space $\R^n$. Dimensions must match exactly.
\end{formaldef}

\begin{formaldef}[Positional Encoding]
The positional encoding matrix $\mP \in \R^{T_{\max} \times d}$ is defined by:
\[
P_{(\text{pos}, 2i)} = \sin\!\left(\frac{\text{pos}}{10000^{2i/d}}\right), \qquad
P_{(\text{pos}, 2i+1)} = \cos\!\left(\frac{\text{pos}}{10000^{2i/d}}\right)
\]
for position $\text{pos} \in \{0, \ldots, T_{\max}-1\}$ and dimension index $i \in \{0, \ldots, d/2 - 1\}$.

The combined input is:
\[
\mX = \mX_{\text{tok}} + \mP_{[:T,:]} \in \R^{T \times d}
\]
\end{formaldef}

\begin{intuition}[Why Addition Works]
Addition blends two signals into one vector. Imagine two radio stations broadcasting
on different frequencies --- you can receive both simultaneously because their signals
don't interfere destructively. Similarly, the learned token embedding and the fixed
sin/cos positional encoding occupy different ``frequency bands'' in the 64 dimensions.

After addition, dimension 0 might be 80\% character information and 20\% position.
Dimension 47 might be the reverse. The subsequent projection matrices (W$_Q$, W$_K$, W$_V$)
can learn to separate or combine these signals as needed.
\end{intuition}

\begin{dimcheck}
\[
\underbrace{\mX_{\text{tok}}}_{T \times d} + \underbrace{\mP}_{T \times d} = \underbrace{\mX}_{T \times d}
\]
Both operands must have identical shape. The output has the same shape.
Addition does not change dimensionality.
\end{dimcheck}

% ============================================================================
\section{Linear Transformations: The Heart of Everything}
% ============================================================================

\begin{formaldef}[Linear Transformation]
A \textbf{linear transformation} (or linear map) is a function $f: \R^n \to \R^m$
that satisfies:
\begin{align}
f(\vx + \bm{y}) &= f(\vx) + f(\bm{y}) & \text{(additivity)} \\
f(\alpha \vx) &= \alpha f(\vx) & \text{(homogeneity)}
\end{align}
Every linear transformation can be represented as multiplication by a matrix
$\mW \in \R^{n \times m}$:
\[
f(\vx) = \vx \mW \qquad \text{where } \vx \in \R^{1 \times n}, \ f(\vx) \in \R^{1 \times m}
\]
The matrix $\mW$ completely defines the transformation. Its columns are the images
of the standard basis vectors.
\end{formaldef}

\begin{intuition}[What a Matrix Does to a Vector]
Multiplying a vector by a matrix does three things geometrically:
\begin{enumerate}[nosep]
  \item \textbf{Rotates} --- changes the direction the vector points
  \item \textbf{Stretches/compresses} --- changes the magnitude along different axes
  \item \textbf{Projects} --- can reduce or increase dimensionality
\end{enumerate}
A $(64, 32)$ matrix takes a 64-dim vector and produces a 32-dim vector.
It ``squeezes'' the information from 64 dimensions into 32, keeping whatever
the matrix has learned is most important. Information is necessarily lost
(64 dimensions into 32 cannot be invertible), but training iterations ensures the
\textit{useful} information is preserved.
\end{intuition}

\begin{formaldef}[Matrix Multiplication: The Dimension Rule]
For matrices $\bm{A} \in \R^{m \times n}$ and $\bm{B} \in \R^{n \times p}$:
\[
\bm{C} = \bm{A} \bm{B} \in \R^{m \times p}, \qquad C_{ij} = \sum_{k=1}^{n} A_{ik} B_{kj}
\]
\textbf{The critical rule}: the inner dimensions must match ($n = n$).
The outer dimensions ($m$ and $p$) survive to become the output shape.
\[
\underbrace{\bm{A}}_{m \times \mathbf{n}} \times \underbrace{\bm{B}}_{\mathbf{n} \times p} = \underbrace{\bm{C}}_{m \times p}
\]
If the inner dimensions don't match, the operation is \textbf{undefined}.
\end{formaldef}

\begin{intuition}[Three Ways to See Matrix Multiplication]
\begin{enumerate}[nosep]
  \item \textbf{Row-column dot products}: entry $C_{ij}$ is the dot product of
        row $i$ of $\bm{A}$ with column $j$ of $\bm{B}$.
  \item \textbf{Transformation of rows}: each row of $\bm{A}$ is a vector that
        gets transformed by $\bm{B}$.
  \item \textbf{Linear combination of columns}: each column of $\bm{C}$ is a
        linear combination of columns of $\bm{A}$, with weights from $\bm{B}$.
\end{enumerate}
In transformers, interpretation (2) is most useful: each position's vector (a row of $\mX$)
is independently transformed by the weight matrix.
\end{intuition}

% ============================================================================
\section{The Q, K, V Projections}
% ============================================================================

\begin{formaldef}[Three Learned Projections]
Given input $\mX \in \R^{T \times d}$ and three weight matrices:
\[
\mW_Q \in \R^{d \times d}, \qquad \mW_K \in \R^{d \times d}, \qquad \mW_V \in \R^{d \times d}
\]
the Query, Key, and Value matrices are:
\[
\mQ = \mX \mW_Q \in \R^{T \times d}, \qquad
\mK = \mX \mW_K \in \R^{T \times d}, \qquad
\mV = \mX \mW_V \in \R^{T \times d}
\]
\end{formaldef}

\begin{dimcheck}
\[
\underbrace{\mX}_{T \times \mathbf{d}} \times \underbrace{\mW_Q}_{\mathbf{d} \times d} = \underbrace{\mQ}_{T \times d}
\]
Inner dimension $d=64$ is consumed. Each of $T$ positions produces a new $d$-dimensional vector.
The same input $\mX$ is multiplied by three \textit{different} matrices to produce three
\textit{different} outputs.
\end{dimcheck}

\begin{intuition}[What Each Projection Does Geometrically]
Each weight matrix rotates and stretches the 64-dimensional input into a new
64-dimensional space:
\begin{itemize}[nosep]
  \item $\mW_Q$ rotates into ``question space'' --- dimensions optimized for expressing what information is needed
  \item $\mW_K$ rotates into ``key space'' --- dimensions optimized for advertising what information is available
  \item $\mW_V$ rotates into ``value space'' --- dimensions optimized for carrying useful content
\end{itemize}
The same input vector $\vx_i$ is viewed through three different ``lenses.''
Each lens extracts different aspects of the combined character+position information.
\end{intuition}

\begin{transformer}[In the Proust Machine]
\texttt{attention.py} lines 199--201:
\begin{verbatim}
Q = x @ self.W_Q   # (1, 13, 64) @ (64, 64) = (1, 13, 64)
K = x @ self.W_K   # (1, 13, 64) @ (64, 64) = (1, 13, 64)
V = x @ self.W_V   # (1, 13, 64) @ (64, 64) = (1, 13, 64)
\end{verbatim}
Three separate applications of the same operation (matrix multiply) with three
different learned matrices. The batch dimension $B$ passes through unchanged ---
each sequence in the batch is processed independently.
\end{transformer}

% ============================================================================
\section{The Dot Product: Measuring Compatibility}
% ============================================================================

\begin{formaldef}[Dot Product (Inner Product)]
For vectors $\vx, \bm{y} \in \R^n$:
\[
\langle \vx, \bm{y} \rangle = \vx^\top \bm{y} = \sum_{i=1}^{n} x_i y_i \in \R
\]
Geometric form:
\[
\langle \vx, \bm{y} \rangle = \|\vx\| \cdot \|\bm{y}\| \cdot \cos\theta
\]
where $\theta$ is the angle between the vectors and $\|\vx\| = \sqrt{\sum x_i^2}$ is the Euclidean norm.
\end{formaldef}

\begin{intuition}[What the Dot Product Means]
The dot product combines magnitude and alignment into one number:
\begin{center}
\begin{tabular}{@{}ll@{}}
\toprule
\textbf{Geometric Relationship} & \textbf{Dot Product} \\
\midrule
Same direction ($\theta = 0$) & Large positive \\
Perpendicular ($\theta = 90°$) & Zero \\
Opposite direction ($\theta = 180°$) & Large negative \\
\bottomrule
\end{tabular}
\end{center}
The dot product does not inherently measure ``meaning similarity.''
It measures geometric alignment. In the transformer, the Q and K projections
\textit{create} a space where alignment corresponds to usefulness --- this
correspondence is "learned" by testing, not built in.
\end{intuition}

\begin{formaldef}[All Pairwise Dot Products via Matrix Multiplication]
To compute the dot product between every pair of rows from $\mQ \in \R^{T \times d_k}$
and $\mK \in \R^{T \times d_k}$:
\[
\mS = \mQ \mK^\top \in \R^{T \times T}, \qquad S_{ij} = \langle \vq_i, \vk_j \rangle = \sum_{l=1}^{d_k} Q_{il} K_{jl}
\]
Entry $S_{ij}$ is the dot product of Query at position $i$ with Key at position $j$:
how compatible position $i$'s question is with position $j$'s advertisement.
\end{formaldef}

\begin{dimcheck}
\[
\underbrace{\mQ}_{T \times \mathbf{d_k}} \times \underbrace{\mK^\top}_{\mathbf{d_k} \times T} = \underbrace{\mS}_{T \times T}
\]
The feature dimension $d_k$ is consumed. What remains is a $T \times T$ matrix:
every position scored against every other position. This is why attention has
$O(T^2)$ complexity --- the score matrix grows quadratically with sequence length.
\end{dimcheck}

\begin{transformer}[In the Proust Machine]
\texttt{attention.py} line 221:
\begin{verbatim}
scores = Q @ K.transpose(0, 1, 3, 2)
# (1, 2, 13, 32) @ (1, 2, 32, 13) = (1, 2, 13, 13)
\end{verbatim}
The transpose swaps the last two dimensions of $\mK$, turning each head's
$(T, d_k)$ into $(d_k, T)$. The batch and head dimensions pass through unchanged.
The result is a $(13, 13)$ compatibility matrix \textit{per head}.
\end{transformer}

% ============================================================================
\section{Scaling: Variance Control}
% ============================================================================

\begin{formaldef}[Scaled Dot-Product]
\[
\mS_{\text{scaled}} = \frac{\mQ \mK^\top}{\sqrt{d_k}}
\]
If entries of $\mQ$ and $\mK$ have variance $\sigma^2$, then entries of $\mQ\mK^\top$
have variance $d_k \cdot \sigma^4$ (sum of $d_k$ products). Dividing by $\sqrt{d_k}$
restores variance to $\sigma^4$, independent of dimension.
\end{formaldef}

\begin{intuition}[Why $\sqrt{d_k}$?]
Without scaling, larger $d_k$ means larger dot products, which means softmax
saturates (outputs near 0 or 1), which means gradients vanish, which means
the model stops learning. The $\sqrt{d_k}$ is not a hyperparameter to tune ---
it's a mathematical correction derived from the statistics of random dot products.

In our model: $\sqrt{d_k} = \sqrt{32} \approx 5.66$. Every score is divided by 5.66.
\end{intuition}

% ============================================================================
\section{The Causal Mask: Constraining the Score Matrix}
% ============================================================================

\begin{formaldef}[Causal Mask]
Define the mask matrix $\bm{M} \in \R^{T \times T}$:
\[
M_{ij} = \begin{cases}
0 & \text{if } j \leq i \quad \text{(past and present)} \\
-\infty & \text{if } j > i \quad \text{(future)}
\end{cases}
\]
Applied by addition: $\mS_{\text{masked}} = \mS_{\text{scaled}} + \bm{M}$
\end{formaldef}

\begin{intuition}[Why Addition of $-\infty$ Works]
Adding $-\infty$ to a score means $\exp(-\infty) = 0$ in the subsequent softmax.
The future position gets exactly zero attention weight. This is more elegant than
deleting entries --- the matrix stays the same shape, and the operation is differentiable
everywhere except at the masked positions (where we don't need gradients anyway).

The mask is a lower-triangular matrix of zeros with $-\infty$ above the diagonal:
\[
\bm{M} = \begin{pmatrix}
0 & -\infty & -\infty & -\infty \\
0 & 0 & -\infty & -\infty \\
0 & 0 & 0 & -\infty \\
0 & 0 & 0 & 0
\end{pmatrix}
\]
\end{intuition}

% ============================================================================
\section{Softmax: From Scores to Probability Distribution}
% ============================================================================

\begin{formaldef}[Softmax Function]
For a vector $\bm{z} \in \R^n$, the softmax function $\sigma: \R^n \to \Delta^{n-1}$
maps to the probability simplex:
\[
\sigma(\bm{z})_i = \frac{e^{z_i}}{\sum_{j=1}^{n} e^{z_j}}, \qquad i = 1, \ldots, n
\]
Properties:
\begin{itemize}[nosep]
  \item $\sigma(\bm{z})_i > 0$ for all $i$ (always positive)
  \item $\sum_i \sigma(\bm{z})_i = 1$ (sums to one)
  \item Monotone: if $z_i > z_j$ then $\sigma(\bm{z})_i > \sigma(\bm{z})_j$ (preserves order)
  \item $\sigma(\bm{z} + c\bm{1}) = \sigma(\bm{z})$ for any constant $c$ (shift-invariant)
\end{itemize}
\end{formaldef}

\begin{formaldef}[Numerically Stable Implementation]
\[
\sigma(\bm{z})_i = \frac{e^{z_i - \max(\bm{z})}}{\sum_{j} e^{z_j - \max(\bm{z})}}
\]
This is mathematically identical (by shift-invariance) but prevents overflow,
since the largest exponent is $e^0 = 1$.
\end{formaldef}

\begin{intuition}[Softmax as Maximum-Entropy Allocation]
Softmax is not an arbitrary normalization. It is the \textbf{unique} probability distribution
that:
\begin{enumerate}[nosep]
  \item Respects the ordering of scores (higher score $=$ higher probability)
  \item Maximizes entropy (makes the fewest assumptions beyond what scores dictate)
\end{enumerate}
This is a theorem from information theory. Any other normalization would inject
unjustified bias into how attention weight is allocated.
\end{intuition}

\begin{transformer}[Applied Row-wise to the Score Matrix]
\[
\mA = \text{softmax}\!\left(\frac{\mQ\mK^\top}{\sqrt{d_k}} + \bm{M}\right) \in \R^{T \times T}
\]
Each row $\mA_{i,:}$ is an independent probability distribution over context positions.
Row $i$ sums to 1 and represents: ``how much attention does position $i$ pay to each
previous position.''

\begin{verbatim}
attention_weights = softmax(scores, axis=-1)
# (1, 2, 13, 13) -- each row of the last two dims sums to 1
\end{verbatim}
\end{transformer}

% ============================================================================
\section{Weighted Sum: Retrieving Information}
% ============================================================================

\begin{formaldef}[Convex Combination of Vectors]
A \textbf{convex combination} of vectors $\vv_1, \ldots, \vv_T \in \R^{d_k}$
with weights $\alpha_1, \ldots, \alpha_T \geq 0$ summing to 1 is:
\[
\vo = \sum_{j=1}^{T} \alpha_j \vv_j = \alpha_1 \vv_1 + \alpha_2 \vv_2 + \cdots + \alpha_T \vv_T \in \R^{d_k}
\]
The result lies within the \textbf{convex hull} of the input vectors --- it cannot
escape the region spanned by them. It is a ``blend'' that lives somewhere
``in between'' the input vectors, positioned by the weights.
\end{formaldef}

\begin{formaldef}[As Matrix Multiplication]
Computing the convex combination for \textit{all} positions simultaneously:
\[
\mO = \mA \mV \in \R^{T \times d_k}
\]
Row $i$ of $\mO$ is the convex combination of all rows of $\mV$, weighted by row $i$ of $\mA$:
\[
\vo_i = \sum_{j=1}^{T} A_{ij} \vv_j
\]
\end{formaldef}

\begin{dimcheck}
\[
\underbrace{\mA}_{T \times \mathbf{T}} \times \underbrace{\mV}_{\mathbf{T} \times d_k} = \underbrace{\mO}_{T \times d_k}
\]
The sequence dimension $T$ is consumed. Each position's output is a $d_k$-dimensional
vector that blends information from all attended positions.
\end{dimcheck}

\begin{intuition}[Soft Selection]
If the attention weights were one-hot (one entry is 1, rest are 0), this would be
\textit{hard} selection: just copy one Value vector. Softmax makes it \textit{soft}:
a weighted average. This softness is critical --- it makes the operation differentiable,
which means gradients can flow through and the model can learn \textit{which} positions
to attend to, using gradient.
\end{intuition}

\begin{transformer}[In the Proust Machine]
\texttt{attention.py} line 248:
\begin{verbatim}
head_outputs = attention_weights @ V
# (1, 2, 13, 13) @ (1, 2, 13, 32) = (1, 2, 13, 32)
\end{verbatim}
Each head independently computes its weighted retrieval. The batch and head
dimensions pass through unchanged. Each position's 32-dim output encodes
context gathered according to that head's attention pattern.
\end{transformer}

% ============================================================================
\section{Reshape and Transpose: Multi-Head Mechanics}
% ============================================================================

\begin{formaldef}[Tensor Reshape]
A \textbf{reshape} reinterprets the memory layout without changing the data.
For a matrix $\bm{A} \in \R^{T \times d}$ where $d = n_h \cdot d_k$:
\[
\text{reshape}(\bm{A}, [T, n_h, d_k]) \in \R^{T \times n_h \times d_k}
\]
The total number of elements is unchanged: $T \cdot d = T \cdot n_h \cdot d_k$.
No computation occurs --- only the ``shape label'' on the data changes.
\end{formaldef}

\begin{formaldef}[Transpose (Axis Permutation)]
A \textbf{transpose} reorders the axes of a tensor. For a 3D tensor $\bm{A} \in \R^{T \times n_h \times d_k}$:
\[
\text{transpose}(\bm{A}, [1, 0, 2]) \in \R^{n_h \times T \times d_k}
\]
This moves the head axis to the front, so we can compute attention for all heads
in a single batched matrix multiplication.
\end{formaldef}

\begin{dimcheck}
Split heads:
\[
\underbrace{\mQ}_{(B, T, d)} \xrightarrow{\text{reshape}} \underbrace{\mQ}_{(B, T, n_h, d_k)} \xrightarrow{\text{transpose}} \underbrace{\mQ}_{(B, n_h, T, d_k)}
\]
In our model:
\[
(1, 13, 64) \to (1, 13, 2, 32) \to (1, 2, 13, 32)
\]

Merge heads (reverse):
\[
\underbrace{\mO}_{(B, n_h, T, d_k)} \xrightarrow{\text{transpose}} \underbrace{\mO}_{(B, T, n_h, d_k)} \xrightarrow{\text{reshape}} \underbrace{\mO}_{(B, T, d)}
\]
\[
(1, 2, 13, 32) \to (1, 13, 2, 32) \to (1, 13, 64)
\]
\end{dimcheck}

\begin{intuition}[Why Move Heads Before Sequence?]
After transposing to $(B, n_h, T, d_k)$, the last two dimensions are $(T, d_k)$ ---
exactly the shape needed for attention: $T$ positions, each with $d_k$ features.
The batch and head dimensions sit in front, meaning NumPy/PyTorch can process
all heads in all batches simultaneously with a single matrix multiplication call.

Without the transpose, we'd need a loop over heads. With it, one \texttt{@} does everything.
\end{intuition}

% ============================================================================
\section{The Output Projection: Mixing Across Heads}
% ============================================================================

\begin{formaldef}[Output Projection]
After merging heads, the concatenated output $\mO_{\text{cat}} \in \R^{T \times d}$
(where the first $d_k$ dimensions are head 1, the next $d_k$ are head 2) is projected:
\[
\mO_{\text{final}} = \mO_{\text{cat}} \mW_O, \qquad \mW_O \in \R^{d \times d}
\]
\end{formaldef}

\begin{dimcheck}
\[
\underbrace{\mO_{\text{cat}}}_{T \times \mathbf{d}} \times \underbrace{\mW_O}_{\mathbf{d} \times d} = \underbrace{\mO_{\text{final}}}_{T \times d}
\]
\end{dimcheck}

\begin{intuition}[Why This Projection Is Necessary]
Before $\mW_O$, the 64 dimensions are rigidly split: dims 0--31 are head 1,
dims 32--63 are head 2. They can't interact. $\mW_O$ is a full $(64, 64)$ matrix
that creates \textit{linear combinations across all 64 dimensions}, allowing
head 1's findings to modulate how head 2's findings are interpreted.

Without $\mW_O$, multi-head attention would just be two independent single-head
attentions glued together. With $\mW_O$, the heads become a coordinated system.
\end{intuition}

% ============================================================================
\section{The Complete Attention Formula}
% ============================================================================

\begin{formaldef}[Single-Head Attention]
\[
\text{Attention}(\mQ, \mK, \mV) = \underbrace{\text{softmax}\!\left(\frac{\mQ\mK^\top}{\sqrt{d_k}} + \bm{M}\right)}_{\mA \in \R^{T \times T}} \mV \in \R^{T \times d_k}
\]
\end{formaldef}

\begin{formaldef}[Multi-Head Attention]
\begin{align}
\text{head}_h &= \text{Attention}(\mX\mW_Q^{(h)},\ \mX\mW_K^{(h)},\ \mX\mW_V^{(h)}) \in \R^{T \times d_k} \\[6pt]
\text{MultiHead}(\mX) &= \text{Concat}(\text{head}_1, \ldots, \text{head}_{n_h}) \, \mW_O \in \R^{T \times d}
\end{align}
\end{formaldef}

\begin{formaldef}[Complete Dimension Flow]
\[
\boxed{
\begin{aligned}
\mX &\in \R^{T \times d} & & \text{input} \\
\mQ = \mX \mW_Q &\in \R^{T \times d} & & \text{projection} \\
\mQ \to \mQ^{(h)} &\in \R^{n_h \times T \times d_k} & & \text{split heads} \\
\mQ^{(h)} {\mK^{(h)}}^\top &\in \R^{n_h \times T \times T} & & \text{score matrix} \\
\frac{1}{\sqrt{d_k}} \cdot (\cdot) + \bm{M} &\in \R^{n_h \times T \times T} & & \text{scale + mask} \\
\text{softmax}(\cdot) = \mA &\in \R^{n_h \times T \times T} & & \text{attention weights} \\
\mA \mV^{(h)} &\in \R^{n_h \times T \times d_k} & & \text{weighted retrieval} \\
\text{merge, project by } \mW_O &\in \R^{T \times d} & & \text{output}
\end{aligned}
}
\]
Input shape equals output shape: $\R^{T \times d} \to \R^{T \times d}$.
\end{formaldef}

% ============================================================================
\section{Summary: Every Operation at a Glance}
% ============================================================================

\begin{center}
\renewcommand{\arraystretch}{1.4}
\begin{tabular}{@{}p{3cm}p{3.5cm}p{3cm}p{4cm}@{}}
\toprule
\textbf{Operation} & \textbf{Linear Algebra} & \textbf{Dimensions} & \textbf{What It Does} \\
\midrule
Embedding lookup & Row selection ($\bm{H}\mE$) & $(T,V) \times (V,d) = (T,d)$ & Discrete symbols to continuous vectors \\
\addlinespace
Position addition & Vector addition & $(T,d) + (T,d) = (T,d)$ & Injects position information \\
\addlinespace
Q, K, V projection & Matrix multiply & $(T,d) \times (d,d) = (T,d)$ & Rotates into Q/K/V spaces \\
\addlinespace
Head split & Reshape + transpose & $(T,d) \to (n_h,T,d_k)$ & Separates heads for parallel attention \\
\addlinespace
Score computation & $\mQ\mK^\top$ & $(T,d_k) \times (d_k,T) = (T,T)$ & All pairwise compatibilities \\
\addlinespace
Scaling & Scalar division & $(T,T) / \sqrt{d_k} = (T,T)$ & Variance normalization \\
\addlinespace
Masking & Matrix addition & $(T,T) + (T,T) = (T,T)$ & Blocks future positions \\
\addlinespace
Softmax & $\R^T \to \Delta^{T-1}$ (per row) & $(T,T) \to (T,T)$ & Scores to probability distributions \\
\addlinespace
Weighted sum & $\mA\mV$ & $(T,T) \times (T,d_k) = (T,d_k)$ & Context-weighted information retrieval \\
\addlinespace
Head merge & Transpose + reshape & $(n_h,T,d_k) \to (T,d)$ & Concatenate head outputs \\
\addlinespace
Output projection & Matrix multiply & $(T,d) \times (d,d) = (T,d)$ & Mixes information across heads \\
\bottomrule
\end{tabular}
\end{center}

% ============================================================================
\section{Key Identities and Properties}
% ============================================================================

\begin{formaldef}[Properties Used in Attention]
\begin{enumerate}[nosep]
  \item \textbf{Associativity}: $(\bm{A}\bm{B})\bm{C} = \bm{A}(\bm{B}\bm{C})$
        --- we can choose the most efficient multiplication order.

  \item \textbf{Transpose of product}: $(\bm{A}\bm{B})^\top = \bm{B}^\top \bm{A}^\top$
        --- order reverses under transpose.

  \item \textbf{Distributivity}: $\bm{A}(\bm{B} + \bm{C}) = \bm{A}\bm{B} + \bm{A}\bm{C}$
        --- used implicitly when the residual connection adds back the input.

  \item \textbf{Dot product as matrix multiply}: $\langle \vx, \bm{y} \rangle = \vx^\top \bm{y}$
        --- a special case of matrix multiplication with $(1, n) \times (n, 1) = (1, 1)$.

  \item \textbf{Batch broadcasting}: a $(d, d)$ weight matrix applied to $(B, T, d)$
        acts independently on each of the $B \times T$ vectors. The batch and sequence
        dimensions are ``spectators'' that pass through unchanged.

  \item \textbf{Softmax shift-invariance}: $\sigma(\bm{z} + c) = \sigma(\bm{z})$
        --- enables the numerical stability trick.

  \item \textbf{Convex combination stays in convex hull}: the output of $\mA\mV$
        (with $\mA$ being row-stochastic) lies within the convex hull of $\mV$'s rows.
        Attention cannot create new information, only \textit{blend} existing information.
\end{enumerate}
\end{formaldef}

% ============================================================================
\section{What This Document Does Not Cover (Yet)}
% ============================================================================

The following components also use linear algebra and will be documented
as we implement them:

\begin{itemize}[nosep]
  \item \textbf{Layer Normalization}: normalization across the feature dimension,
        using mean and variance statistics. Involves element-wise operations and
        learned affine parameters.
  \item \textbf{Feedforward Network}: two linear transformations with a nonlinearity:
        $\text{FFN}(\vx) = \mW_2 \, \text{GELU}(\mW_1 \vx + \bm{b}_1) + \bm{b}_2$.
        Changes dimension $d \to 4d \to d$.
  \item \textbf{Residual Connection}: $\vx + f(\vx)$, simple vector addition
        that enables gradient flow through deep networks.
  \item \textbf{Final Linear Projection}: $(T, d) \times (d, V) = (T, V)$,
        mapping from model space to vocabulary logits.
  \item \textbf{Backpropagation}: how gradients of the loss with respect to
        each weight matrix are computed using the chain rule and transposed Jacobians.
\end{itemize}

\vfill
\begin{center}
\textit{``The only real voyage of discovery consists not in seeking new landscapes,\\
but in having new eyes.'' --- Marcel Proust}
\end{center}

\end{document}
